\begin{table}[t]
\centering
\caption{Standardized Effect Sizes}
\label{tab:sde}
\small
\begin{adjustbox}{max width=\textwidth}
\begin{tabular}{lcccccc}
\hline\hline
Outcome & $\hat{\beta}$ & SE & SD($Y$) & SDE & SE(SDE) & Classification \\
\hline
\multicolumn{7}{l}{\textit{Panel A: Pooled}} \\
Enrollment volatility & 0.627 & 0.173 & 1.006 & 0.274 & 0.076 & Large positive \\
Rolling CV (12-month) & 0.037 & 0.022 & 0.020 & 0.838 & 0.484 & Large positive \\
Log applications & 0.725 & 0.314 & 1.296 & 0.246 & 0.107 & Large positive \\
\addlinespace
\multicolumn{7}{l}{\textit{Panel B: Heterogeneous (by Medicaid expansion status)}} \\
Enrollment vol. (expansion) & 0.731 & 0.699 & 1.081 & 0.298 & 0.285 & Large positive \\
Enrollment vol. (non-expansion) & 2.662 & 1.687 & 0.833 & 1.407 & 0.892 & Large positive \\
\hline\hline
\end{tabular}
\end{adjustbox}
\begin{tablenotes}
\item \textit{Notes:} \textbf{Country:} United States. \textbf{Research question:} Does SNAP recertification intensity increase Medicaid enrollment volatility in states with integrated eligibility systems? \textbf{Policy mechanism:} Integrated eligibility systems process SNAP and Medicaid through shared caseworkers and IT infrastructure; higher SNAP recertification frequency creates administrative backlogs that spill over to Medicaid enrollment processing. \textbf{Outcome definition:} Absolute month-over-month percent change in total state Medicaid enrollment (primary); 12-month rolling coefficient of variation (secondary); log new Medicaid applications (mechanism). \textbf{Treatment:} Continuous --- share of SNAP households on certification periods of six months or less, interacted with binary indicator for integrated eligibility system status. \textbf{Data:} CMS Monthly Medicaid Enrollment Reports and USDA ERS SNAP Policy Database, January 2018 to December 2020, state-month panel, 1,785 observations. \textbf{Method:} OLS with state and year-month fixed effects, standard errors clustered at state level. \textbf{Sample:} 51 U.S.\ states including DC; 26 integrated eligibility system states, 25 non-integrated. SDE $= \hat{\beta} \times \text{SD}(X) / \text{SD}(Y)$ where SD($X$) is the standard deviation of recertification intensity and SD($Y$) is the sample standard deviation of the outcome. Classification refers to magnitude, not statistical significance: Large ($|$SDE$| > 0.15$), Moderate ($0.05$--$0.15$), Small ($0.005$--$0.05$), Null ($< 0.005$).
\end{tablenotes}
\end{table}

